\documentclass{beamer}
%\usepackage{minted}
%\usepackage{epigraph}
\usepackage{subfig}
\usepackage{pgfplots}
\pgfplotsset{compat=1.18}
\usepackage{tikz}
\usepackage{color, colortbl}
\usepackage{amsmath}
\usepackage{tabularray}
\usepackage{verbatim}
\usepackage[
    n,
    operators,
    advantage,
    sets,
    adversary,
    landau,
    probability,
    notions,    
    logic,
    ff,
    mm,
    primitives,
    events,
    complexity,
    asymptotics,
    keys]{cryptocode}
\usepackage{bm}
\usepackage{charter}
\usepackage{environ}
\usepackage{tikz}
\usetikzlibrary{calc,matrix}

%% Code by Claudio:
%% https://tex.stackexchange.com/a/197447/221452
%% Uses code by Andrew:
%% http://tex.stackexchange.com/a/28452/13304
\makeatletter
    \let\matamp=&
    \catcode`\&=13
    \def&{%
        \iftikz@is@matrix%
            \pgfmatrixnextcell%
        \else%
            \matamp%
        \fi%
    }
\makeatother

\newcounter{lines}
\def\endlr{\stepcounter{lines}\\}

\newcounter{vtml}
\setcounter{vtml}{0}

\newif\ifvtimelinetitle
\newif\ifvtimebottomline

\tikzset{
    description/.style={column 2/.append style={#1}},
    timeline color/.store in=\vtmlcolor,
    timeline color=red!80!black,
    timeline color st/.style={fill=\vtmlcolor,draw=\vtmlcolor},
    use timeline header/.is if=vtimelinetitle,
    use timeline header=false,
    add bottom line/.is if=vtimebottomline,
    add bottom line=false,
    timeline title/.store in=\vtimelinetitle,
    timeline title={},
    line offset/.store in=\lineoffset,
    line offset=4pt,
}

\NewEnviron{vtimeline}[1][]{%
    \setcounter{lines}{1}%
    \stepcounter{vtml}%
    \begin{tikzpicture}[column 1/.style={anchor=east},
        column 2/.style={anchor=west},
        text depth=0pt,text height=1ex,
        row sep=1ex,
        column sep=1em,
        #1
    ]
        \matrix(vtimeline\thevtml)[matrix of nodes]{\BODY};
        \pgfmathtruncatemacro\endmtx{\thelines-1}

        \path[timeline color st]
            ($(vtimeline\thevtml-1-1.north east)!0.5!(vtimeline\thevtml-1-2.north west)$)--
            ($(vtimeline\thevtml-\endmtx-1.south east)!0.5!(vtimeline\thevtml-\endmtx-2.south west)$);

        \foreach \x in {1,...,\endmtx}{
            \node[circle,timeline color st, inner sep=0.15pt, draw=white, thick]
            (vtimeline\thevtml-c-\x) at
            ($(vtimeline\thevtml-\x-1.east)!0.5!(vtimeline\thevtml-\x-2.west)$){};
                \draw[timeline color st](vtimeline\thevtml-c-\x.west)--++(-3pt,0);
        }

        \ifvtimelinetitle%
            \draw[timeline color st]([yshift=\lineoffset]vtimeline\thevtml.north west)--
                ([yshift=\lineoffset]vtimeline\thevtml.north east);

            \node[anchor=west,yshift=16pt,font=\large]
                at (vtimeline\thevtml-1-1.north west)
                {\textsc{Timeline \thevtml}: \textit{\vtimelinetitle}};
        \else%
            \relax%
        \fi%

        \ifvtimebottomline%
            \draw[timeline color st]([yshift=-\lineoffset]vtimeline\thevtml.south west)--
            ([yshift=-\lineoffset]vtimeline\thevtml.south east);
        \else%
            \relax%
        \fi%
    \end{tikzpicture}
}
\usetheme{focus}

\title{RUSTico}
\subtitle{Rust in cryptographic outlook}
\author{Davide Carnemolla}
\titlegraphic{\includegraphics[scale=0.12]{images/ferris.png}}
\institute{Università degli Studi di Catania}
\date{2024}

\begin{document}
    \begin{frame}
        \maketitle
    \end{frame}

    \section{Introduction}

    \begin{frame}{Introduction: main features}
        \begin{columns}[t, onlytextwidth]
            \column{0.33\textwidth}
            \centering
            \includegraphics[height=1.5cm, keepaspectratio]{images/shapes.pdf}

            \textbf{General purpose}

            \column{0.33\textwidth}
            \centering
            \includegraphics[height=1.5cm, keepaspectratio]{images/fast.pdf}
            
            \textbf{Blazingly fast}

            \column{0.33\textwidth}
            \centering
            \includegraphics[height=1.5cm, keepaspectratio]{images/memory.pdf}
            
            \textbf{Memory safe}
        \end{columns}

        \vspace{1cm}

        \begin{columns}[t, onlytextwidth]
            \column{0.5\textwidth}
            \centering
            \includegraphics[height=1.5cm, keepaspectratio]{images/no-gc.pdf}

            \textbf{No Garbage Collector}

            \column{0.5\textwidth}
            \centering
            \includegraphics[height=1.5cm, keepaspectratio]{images/toolbox.pdf}

            \textbf{Tools} \\ (doc, cargo, libraries)
        \end{columns}
    \end{frame}

    \begin{frame}{Introduction: history}
        \begin{vtimeline}[description={text width=7cm},
        row sep=4ex]
        2006 & Mozilla creates Rust for the Servo's project \endlr
        2010 & Rust is announced at Mozilla Summit\endlr
        2012 & Alpha release\endlr
        2013 & Version 0.6 release\endlr
        2015 & First stable release (Rust 1.0)\endlr
        2021 & The Rust Foundation is founded\endlr
    \end{vtimeline}
    \end{frame}

    \begin{frame}{Introduction: who uses it?}
        \begin{itemize}
            \item AWS: Firecracker powers Lambda and Fargate
            \item Google: parts of the Fuschia operating system
            \item Linux: 2nd official language for the Kernel!
            \item CloudFlare: quic / http 3 implementation
            \item Dropbox: file storage
            \item Clever Cloud: reverse proxy
            \item Atlassian, Canonical, Coursera, Chef, Deliveroo, NPM, Sentry...
            \item Growing ecosystem for embedded development
        \end{itemize}
    \end{frame}

    \begin{frame}{Introduction: Rust vs C/C++}
        \begin{center}
            \begin{tabular}{ |c|c|c|c| } 
             \hline
             \textbf{Properties} & \textbf{Rust} & \textbf{C} & \textbf{C++} \\
             \hline
             Type safety & Safe & Unsafe & Unsafe \\
             \hline
             Type expression & Implicit/Explicit & Explicit & Explicit \\
             \hline
             Type checking & Static/dynamic & Static & Static \\
             \hline
             Garbage collector & No & No & No \\
             \hline
            \end{tabular}
            \end{center}
    \end{frame}
    

    \begin{frame}{Introduction: learning Rust}
        \begin{columns}[t, onlytextwidth]
            \column{0.5\textwidth}
            \centering
            \includegraphics[height=5.5cm, keepaspectratio]{images/rust-book.jpg}

            \small Online at https://www.rust-lang.org

            \column{0.5\textwidth}
            \centering
            \includegraphics[height=5.5cm, keepaspectratio]{images/programming-rust.jpg}
        \end{columns}
    \end{frame}

    \begin{frame}[containsverbatim]{Getting started}
        \small
        \begin{block}{Installing on MacOS/Unix}
            \tt{curl --proto '=https' --tlsv1.2 -sSf https://sh.rustup.rs | sh}
        \end{block}

        \begin{block}{Installing on Windows}
            Download and install the \tt{rustup-init.exe} from the official website. 
        \end{block}

        \begin{alertblock}{Verifying the installation}
            \tt{rustc --version}
        \end{alertblock}
    \end{frame}
    
\end{document}
